\documentclass[a5paper,10pt]{article}

% https://github.com/InDevRus/filippov-solutions

% Нумерация страниц
\pagenumbering{gobble}

% Настройка полей
\usepackage{geometry}
\geometry{left=1cm}
\geometry{right=1cm}
\geometry{top=1cm}
\geometry{bottom=1.5cm}

% Вставка таблицы
\usepackage{array}
\usepackage{tabularx}

% Инструменты выравнивания формул
\usepackage{amsmath}
\usepackage{mathtools}

% Диакритика в математике
\usepackage{amsfonts}

% Вычисления размеров на ходу
\usepackage{calc}

% Удобные записи производных
\usepackage[italic=true]{derivative}

% Настройки сносок
\usepackage{enotez}

% Длинные знаки векторов
\usepackage{anyfontsize}
\usepackage[b]{esvect}

% Вставка картинок
\usepackage{graphicx}

% Вставка красной строки
\usepackage{indentfirst}
\setlength{\parindent}{1em}

% Условия в командах
\usepackage{xifthen}

% Луа-скрипты
\usepackage{luacode}

% Настройка команд отдельно в математическом и текстовом режимах
\usepackage{mathcommand}

% Для повторения LaTeX кода
\usepackage{multido}

% Конфигурация отступов формул
\delimitershortfall-1sp
\usepackage{mleftright}
\mleftright

% Растягивание объектов
\usepackage{relsize}

% Обтекание текстом
\usepackage{wrapfig}
\usepackage{subcaption}
\usepackage{floatflt}

% Поддержка ввода кириллицы
\usepackage[english, russian]{babel}

% Настройка корневой позиции для компилятора
\usepackage{import}
\providecommand*{\ProjectRootPath}{..}

\import{\ProjectRootPath/preambles/macros/}{theorems}

\import{\ProjectRootPath/preambles/fonts}{font_setup}

% Знак градуса
\usepackage{siunitx}

\import{\ProjectRootPath/preambles/macros/}{operators}
\import{\ProjectRootPath/preambles/macros/}{redefinitions}

\import{\ProjectRootPath/preambles/fonts}{letters_setup}
\import{\ProjectRootPath/preambles/fonts/adjustments}{custom_superscripts}

\import{\ProjectRootPath/preambles/custom}{formatting}
\import{\ProjectRootPath/preambles/custom}{frames}
\import{\ProjectRootPath/preambles/custom}{list_setup}

% TODO: перевести команды на современный стиль объявления

\begin{document}

\begin{framed}
    Если в полярных координатах
    $\tsys{x = r \cos\theta \\ y = r \sin\theta}$
    имеем производную 
    $\der{r}{\theta} = A\br{r; \theta}$,
    то соответствующая производная в декартовых координатах 
    $\der{y}{x} = B\br{r; \theta}$ связана с исходной $A\br{r; \theta}$ соотношениями
    \begin{align*}
        \der{y}{x}
        = B\br{A; r; \theta}
        = \dfrac {A \sin\theta + r \cos\theta} {A \cos\theta - r \sin\theta};&&
        \der{r}{\theta}
        = A\br{B; r; \theta}
        = r \cdot \dfrac {B \sin\theta + \cos\theta} {B\cos\theta - \sin\theta}.
    \end{align*}
\end{framed}

\begin{framed}
    Для угла $\phi = \ang{90}$ имеем 
    $\tg{\beta} 
    = \tg\br{\alpha \pm \phi} 
    = \tg\br{\alpha \pm \ang{90}} 
    = -\ctg \alpha 
    = -\dfrac {1} {\tg\alpha}$.
    
    С учётом указанного выше для данного угла $\phi$ 
    всякое семейство изогональных траекторий $z\br{x}$ 
    и соответствующее семейство в полярной системе 
    $\tsys{x = \rho \cos\theta \\ z = \rho \sin\theta}$
    удовлетворяют соотношениям
    \begin{align*}
        \sub{A}{z}
        = \rho \cdot \dfrac 
        {\sub{B}{z} \sin\theta + \cos\theta} 
        {\sub{B}{z}\cos\theta - \sin\theta};
        &&
        \sub{B}{z}\br{\rho; \theta} 
        = -\dfrac {1} {\sub{B}{y}\br{\rho; \theta}};
        &&
        \sub{B}{y} = \dfrac 
        {\sub{A}{y} \sin\theta + r \cos\theta}
        {\sub{A}{y} \cos\theta - r \sin\theta},
    \end{align*}
    где введены обозначения
    $\sub{A}{y} = \der{r}{\theta}$,
    $\sub{B}{y} = \der{y}{x}$,
    $\sub{A}{z} = \der{\rho}{\theta}$,
    $\sub{B}{z} = \der{z}{x}$.
    
    Последовательно подставив и затем упростив, получим следующую формулу:
    $$\sub{A}{z}\br{\rho; \theta} 
    = -\dfrac {\pow{\rho}{2}} {\sub{A}{y}\br{\rho; \theta}}.$$
\end{framed}

$r = a \cos^{2}\theta$.
$a = \dfrac {r} {\cos^{2} \theta}$.
$0 = \dfrac 
{r\` \cos^{2} \theta + 2r \sin\theta \cos\theta} 
{\cos^{4} \theta}$.
$r\` = -2r \tg\theta = \sub{A}{y}$.
Таким образом, уравнение изогональных (ортогональных) траекторий 
в обозначениях, указанных выше, принимает форму
$\der{\rho}{\theta} 
= \sub{A}{z} 
= -\dfrac {\pow{\rho}{2}} {\sub{A}{y}} 
= \dfrac {\pow{\rho}{2}} {2\rho \tg\theta}.$
$\rho\`\br{\theta} = \dfrac {1} {2} \rho\br{\theta} \ctg\theta$.

\begin{remark}
    Уравнение
    $\rho\` = \dfrac {1} {2} \rho \ctg\theta$
    имеет общее решение $\rho = C \sqrt{\vbr{\sin\theta}}$.
\end{remark}

\end{document}

